%%%%%%%%%%%%%%%%%%%%%%%%%%%%%%%%%%%%%%%%%%%%%%%%%%%%%%%%%%%%%%%%%
\chapter{TARTIŞMA}\label{CH5}
%%%%%%%%%%%%%%%%%%%%%%%%%%%%%%%%%%%%%%%%%%%%%%%%%%%%%%%%%%%%%%%%%

\section{Araştırma Sorularının Yanıtları}\label{sec:yanitlar}

\textbf{AS1: Meta veri.} ROI'nin konumu, boyutu ve şekli tek başına teşhisi büyük ölçüde ele vermektedir: beyin MR verisinde tümör tipi AUC $0.83$, göğüs röntgeninde teşhis AUC $0.83$--$0.90$. Görüntülerin orijinal çözünürlükte gönderildiği bir sistemde girdi boyutu da teşhisi ele verebilmektedir (AUC $0.93$). Kesit yönü kontrolü, beyin MR verisindeki sızıntının yönden kaynaklanmadığını göstermiştir. $\Pi_{\mathrm{ROI}}$ protokolünün sızıntı fonksiyonu bu bilgilerin tamamını sunucuya serbest bırakmaktadır \cite{bakas2026roi}.

\textbf{AS2: Açık bağlam.} ROI dışındaki bağlam teşhisi neredeyse ROI'nin kendisi kadar iyi ele vermektedir. Beyin MR verisinde tümör gizliyken AUC $0.977$, göğüs röntgeninde akciğerler gizliyken $0.993$ bulunmuştur (5 tohum). Başka bir kaynaktan veriyle eğitilen saldırgan da pnömoniyi AUC $0.934$--$0.998$ ile tahmin etmiştir. Bu nedenle tehdit, eğitim ve testin aynı kümeden gelmesinin bir yan etkisi değildir.

\textbf{AS3: Sızıntının kaynağı.} Saldırgan akciğerin dışındaki geniş bir bağlamı kullanmaktadır: görüntü kenarı, kalp ve mediasten, göğüs duvarı ve ROI'nin şekli. Beyin MR görüntüsünde beynin geri kalanı ve kafa dışı arka plan öne çıkmıştır. Histogram eşitleme ve kenar gizleme gibi basit önişleme adımları sızıntıyı gidermemiştir. Bu bulgu, sızıntının yalnızca görüntü kenarındaki yazı ve işaret gibi yüzeysel izlerden kaynaklanmadığını göstermektedir.

\textbf{AS4: Bedel.} Saldırgan AUC'sini $0.8$'in altına indirmek için göğüs röntgeninin \%98'inin, beyin MR görüntüsünün ise tamamının şifrelenmesi gerekmiştir. Bu noktalarda protokolün hız kazancı $512 \times 512$ görüntülerde sırasıyla 1.02 ve 1.00 kattır. Başka bir deyişle, gizlilik şartı konduğu anda seçici şifrelemenin verimlilik avantajı ortadan kalkmaktadır.

\textbf{AS5: Sızıntısız alternatif.} Mümkündür. FoveaHE'de sunucunun gördüğü bilgilerle kurulan saldırgan şans düzeyinde kalmıştır (AUC $0.45$--$0.53$, $p \geq 0.20$). Aynı model sınıfında tam şifrelemeye göre hız kazancı beyin MR görüntüsünde 29--37, göğüs röntgeninde 7--9 kattır ve bir görüntünün şifreli çıkarımı 0.5--1.5 saniye sürmektedir. Teşhis başarısı tam görüntüyle eğitilen aynı model sınıfının üstündedir: beyin MR verisinde $0.881$'e karşı $0.948$, göğüs röntgeninde $0.946$'ya karşı $0.954$. Deneylerden önce belirlenen beş başarı ölçütünün tamamı 5 tohumla karşılanmıştır; tek istisna, göğüs röntgeninde Model C'nin F64\_G32 yapılandırmasında 4.4 katta kalan hız kazancıdır.

\section{Seçici Şifreleme Literatürü İçin Anlamı}\label{sec:anlam}

$\Pi_{\mathrm{ROI}}$ protokolünün güvenlik kanıtı doğrudur: sunucu, sızıntı fonksiyonunun izin verdiğinden fazlasını öğrenmez. Sorun, izin verilen sızıntının tıbbi görüntüler için fazla cömert olmasıdır. Bu durum, şifreli veritabanlarında izin verilen sızıntının istismarını gösteren çalışmalarla aynı yapıdadır \cite{naveed2015inference,cash2015leakage}. Tıbbi görüntülerde ROI dışındaki bölge ile ROI geometrisi, şifreli bölgenin klinik içeriğiyle güçlü biçimde ilişkilidir ve bu nedenle ``hassas olmayan'' kabul edilemez.

``Encrypt What Matters'' çalışması, ROI şifreli çıkarımda hız kazancının ROI boyutu küçüldükçe ve mimarinin yerelliği arttıkça büyüdüğünü göstermiştir \cite{backour2026encrypt}. Bu tezin bulguları, gizlilik şartı konduğunda şifrelenmesi gereken alanın görüntünün neredeyse tamamına ulaştığını göstermektedir. Bu nedenle o çalışmadaki yerellik avantajının da benzer biçimde azalması beklenir. Ancak derin evrişimli ağlar için şifreli çıkarım maliyeti bu tezde ölçülmemiştir.

Tıbbi görüntülerdeki kısayol öğrenme olgusu \cite{degrave2021shortcuts,maguolo2021critic}, tanı modelleri açısından bir genelleme sorunu olarak ele alınmıştır. Bu tez aynı olgunun seçici şifreleme açısından bir gizlilik açığı olduğunu göstermektedir. Sunucu açısından, teşhisi ele veren ipucunun klinik bir bulgu mu yoksa çekim koşullarına bağlı bir iz mi olduğu fark etmez; her iki durumda da şifreli bölgenin içeriği sızmaktadır.

Bu bulgular Bölüm \ref{sec:privacyutility}'teki kuramsal açıklamayla da uyumludur. Sunucunun teşhis hesaplamasında yararlandığı açık bilgi, teşhisle ilişkili olduğu ölçüde saldırgan için de bilgi taşır. Bu nedenle tıbbi görüntülerde açık bırakılabilecek ``zararsız'' bir bölge aramak yerine verimliliği şifrelenen verinin boyutundan kazanmak gerekir. Tezde önerilen FoveaHE bu ilkenin bir uygulamasıdır.

\section{Önerilen Yöntemin Değerlendirmesi}\label{sec:fovtartisma}

\textbf{Neden çalışıyor?} FoveaHE'nin başarısı üç özellikten gelmektedir. Birincisi, temsilin tamamı şifreli olduğu için bağlam saldırısının kullanacağı açık piksel yoktur. İkincisi, temsilin boyutu her görüntüde aynı olduğundan paket boyutu ve şifreli metin sayısı da sabittir; meta veri saldırısının dayandığı ROI geometrisi yalnızca şifreli içerikte kalır. Üçüncüsü, sığ şifreli modeller ROI'yi kendi başlarına bulamazken odaklı temsil ROI'yi modele merkezlenmiş ve yüksek çözünürlüklü olarak sunar. Bu nedenle FoveaHE yalnızca daha hızlı değil, tam görüntüyle eğitilen aynı model sınıfından daha doğrudur.

Bilgi üst sınırı deneyi, beyin MR verisindeki kazancın önemli bölümünün ROI'nin konum bilgisinden geldiğini göstermiştir; ancak odaklama, aynı bilgi verildiğinde de küçük ve tutarlı bir katkı sağlamaktadır. Göğüs röntgeninde akciğer penceresi görüntünün neredeyse tamamı olduğundan bilgi düzeyinde odaklamanın katkısı görülmemiş; buna karşın sığ modellerde ROI'nin merkezlenmesi eş bütçeli küçültmeye göre belirgin kazanç sağlamıştır.

\textbf{Şifreli özetle karşılaştırma.} Şifreli özet yaklaşımı da açık bilgi bırakmaz, hızlıdır ve iki veri kümesinde de FoveaHE'den daha doğrudur (beyin MR verisinde $0.007$, göğüs röntgeninde $0.027$ fark). Bu sonuç açıkça kabul edilmelidir: tezin katkısı en doğru sızıntısız yöntemi önermek değildir. İki yaklaşımın ödünleşimleri farklıdır. Şifreli özet, istemcide 11.2 milyon parametreli, genel amaçlı bir ağın çalıştırılmasını gerektirir ve başarısı bu ağın öğrendiği özniteliklerin göreve uygunluğuna bağlıdır. FoveaHE ise istemcide yalnızca yeniden örnekleme yapar; temsili görüntü biçimindedir, bütçesi ve odak payı açıkça ayarlanabilir ve hangi bölgenin hangi çözünürlükte gönderildiği bellidir. İki yaklaşım birbirinin rakibi olmaktan çok tamamlayıcısıdır: odak ve genel bakışın önceden eğitilmiş bir ağla ayrı ayrı özetlenip birlikte şifrelenmesi doğal bir sonraki adımdır.

\textbf{Bozuk bağlamın başarısızlığı.} Açık bağlamı gürültü veya bulanıklıkla bozmak, en güçlü düzeylerde bile saldırgan başarısını en fazla $0.01$ azaltmıştır. Bu sonuç, görüntü gizleme yöntemlerinin derin öğrenmeyle aşılabildiğini gösteren bulgularla \cite{mcpherson2016obfuscation} ve Bölüm \ref{sec:privacyutility}'teki açıklamayla tutarlıdır: bozulma faydayı koruduğu sürece teşhis izini de korumaktadır. Bu nedenle ``açık bırak ama boz'' yaklaşımı ROI tabanlı seçici şifrelemenin sızıntısına çözüm değildir.

\textbf{Bütçe ve genelleme.} Bütçe duyarlı odak deneyi, sabit bir değer bütçesinde çözünürlüğü odak ile genel bakış arasında paylaştırmanın her iki uç durumdan da iyi olduğunu göstermiştir. Beyin MR verisinde 1.024--2.048 değer yeterlidir; bu, tek bir şifreli metnin kapasitesinin küçük bir bölümüdür. Harici derlemede temsiller arasındaki sıralama korunmuş ve odaklamanın genellemeyi bozmadığı görülmüştür. Harici kümenin daha kolay olması nedeniyle bu sonuç göreli bir karşılaştırma olarak değerlendirilmelidir.

\section{Tasarım Kuralları}\label{sec:kurallar}

Bulgulara dayanarak homomorfik şifreli tıbbi görüntü çıkarımı kullanacak sağlık sistemleri için aşağıdaki kurallar önerilmektedir:

\begin{enumerate}
\item \textbf{ROI geometrisi sunucuya görünmemelidir.} Görüntüye göre değişen ROI, konumu, şekli ve paket boyutuyla teşhisi ele verir. Sunucuya gönderilen verinin boyutu tüm görüntülerde aynı olmalıdır. FoveaHE'de sabit boyutlu temsil bu kuralı kendiliğinden sağlar.
\item \textbf{Teşhis gizlenecekse açık piksel bırakılmamalı, verimlilik çözünürlükten kazanılmalıdır.} Bulgular, gizlilik şartı konduğunda şifreli alanın görüntünün neredeyse tamamına ulaştığını göstermektedir. Bu nedenle verimlilik şifrelenen alanı değil, şifrelenen veriyi küçülterek sağlanmalıdır. FoveaHE ve şifreli özet bu kuralın iki uygulamasıdır.
\item \textbf{Açık bilgi bırakan her tasarımın sızıntısı ölçülmelidir.} Bir sistem yine de açık bilgi bırakacaksa, bırakılacak bilginin kapsamı temsil edici veri üzerinde savunmaya uyum sağlayan bir saldırganla ölçülen sızıntıya dayanmalıdır. Gürültü ekleme veya bulanıklaştırma sızıntıyı gidermez. Bu tezde geliştirilen değerlendirme boru hattı bu amaçla kullanılabilir.
\item \textbf{Çözülmüş sonuç sunucuya geri gönderilmemelidir.} Çözülmüş CKKS sonuçlarının paylaşılması gizli anahtarın kurtarılmasına yol açabilir \cite{limicciancio2021,guo2024keyrecovery}. Bu tez kapsamında yapılan kavram kanıtı denemesinde, TenSEAL ile ($N = 8192$) toplu işlenen iki şifreli sonucun tam hassasiyetli çözülmüş değerleri paylaşıldığında gizli anahtar 0.15 saniyede kurtarılmıştır. Değerlerin 4 ondalığa yuvarlanması bu basit saldırıyı durdurmuş, ancak bu durum güvenliğin kanıtı sayılmamalıdır.
\item \textbf{Değerlendirme hijyeni sağlanmalıdır.} Şans düzeyine yakın sonuçlar ölçülürken 16 bit hassasiyet ve sınıfa göre sıralı test verisi sahte AUC değerleri üretebilmektedir; bu tezde özdeş girdilerde $0.421$ gözlenmiştir. Tahminler 32 bit hassasiyetle ve karışık sırada hesaplanmalı, tamamen gizli bir kontrol görüşü her deneye eklenmelidir. Sızıntı denetiminde kat ortalamalı AUC, etiket permütasyon testi ve sızıntıyı yakalayabildiği bilinen bir pozitif kontrol kullanılmalıdır.
\end{enumerate}

\section{Tablo Verisinde Seçici Şifreleme Üzerine Not}\label{sec:tablo}

TÜBİTAK önerisi aşamasında öznitelik düzeyinde seçici CKKS şifrelemesi, Kaggle'da yayımlanmış bir inme veri kümesi üzerinde lojistik regresyonla ön çalışma olarak denenmiştir. Bu ön çalışmada üç gözlem yapılmıştır.

\begin{itemize}
\item Tam şifreli lojistik regresyon çıkarımı zaten ucuzdur: tek hasta için yaklaşık 15 ms, 1.022 hastanın toplu işlenmesi için yaklaşık 96 ms. Bu nedenle seçici şifrelemenin maliyet avantajı ihmal edilebilir düzeydedir.
\item Açık bırakılan yarı tanımlayıcılar şifreli öznitelikleri sızdırmaktadır: hipertansiyon AUC $0.79$, kalp hastalığı $0.85$, sigara kullanımı $0.57$--$0.59$.
\item Model skorunun sunucuya açılması, şifreli özniteliklerin doğrusal denklemlerle kurtarılmasına izin verir; bu, dikey federatif öğrenmede gösterilmiş bilinen bir saldırıdır \cite{luo2021feature}.
\end{itemize}

Açık özniteliklerden gizli özniteliğin tahmini ise büyük ölçüde imputasyon niteliğindedir \cite{jayaraman2022imputation}. Tablo verisinde bu olgular literatürde bilindiği ve homomorfik şifrelemenin maliyeti küçük olduğu için, tez özgün katkısını şifreleme maliyetinin gerçekten yüksek olduğu tıbbi görüntüler üzerine kurmuştur.

\section{TÜBİTAK 2209-A Başarı Ölçütlerinin Değerlendirilmesi}\label{sec:tubitakolcut}

Öneride iki başarı ölçütü tanımlanmıştı: şifresiz duruma göre saldırı başarısında en az \%20 düşüş ve tam homomorfik şifrelemeye göre işlem süresinde en az \%30 azalma.

\begin{itemize}
\item \textbf{$\Pi_{\mathrm{ROI}}$ ile süre ölçütü.} Varsayılan kullanımda karşılanmaktadır: $512 \times 512$ görüntüde göğüs röntgeni için yaklaşık 4 kat, beyin MR görüntüsü için 24 kat hızlanma elde edilmiştir.
\item \textbf{$\Pi_{\mathrm{ROI}}$ ile saldırı ölçütü.} Aynı varsayılan kullanımda karşılanmamıştır: saldırgan AUC değeri göğüs röntgeninde $0.997$'den $0.993$'e, beyin MR verisinde $0.985$'ten $0.977$'ye inmiştir. Göğüs röntgeninde \%20'lik düşüş ancak görüntünün \%98'i şifrelendiğinde sağlanmış (AUC $0.776$), beyin MR verisinde ise yalnızca tam şifrelemeyle mümkün olmuştur. İki ölçüt aynı anda karşılanamamıştır; saldırı ölçütünün sağlandığı noktada hız kazancı \%2'nin altındadır.
\item \textbf{FoveaHE ile iki ölçüt birlikte.} İki ölçüt aynı anda karşılanmıştır. İşlem süresi tam şifrelemeye göre beyin MR görüntüsünde yaklaşık \%97, göğüs röntgeninde \%86--89 azalmıştır. Sunucunun gördüğü bilgilerle saldırgan AUC değeri şans düzeyine inmiş ($0.53$ ve $0.50$), şifresiz duruma göre ($0.985$ ve $0.997$) yaklaşık \%47--50 düşmüştür.
\end{itemize}

Bu sonuçlar, önerinin hipotezinin ROI tabanlı seçici şifrelemeyle doğrulanamadığını, ancak çözünürlükte seçicilik ilkesine dayanan FoveaHE ile iki ölçütün birlikte karşılanabildiğini göstermektedir. Böylece önerinin temel sorusuna iki yönlü bir yanıt verilmiştir: tıbbi görüntülerde açık bölge bırakan seçici şifreleme saldırıları anlamlı düzeyde azaltırken maliyeti düşük tutamamakta, temsilin tamamını şifreleyen ve verimliliği çözünürlükten kazanan bir tasarım ise bunu başarabilmektedir.

\section{Sınırlılıklar}\label{sec:sinirliliklar}

\begin{itemize}
\item \textbf{Veri kümeleri.} COVID-QU-Ex birden fazla kaynaktan derlenmiştir. Kaggle kümesinin büyük bölümü COVID-QU-Ex içinde yer almaktadır ve iki küme aynı çocuk hastalar kaynağını paylaşmaktadır. Beyin MR kümesi iki hastaneden gelmektedir. Çekim koşullarına bağlı karıştırıcı etkenler mutlak AUC değerlerini yükseltmiş olabilir. Ancak bu tür etkenler gerçek sistemlerde de bulunduğu için bulgunun güvenlik açısından önemini azaltmaz.
\item \textbf{Tohum sayısı.} Meta veri saldırısı, bağlam saldırısının ana görüşleri ve önerilen yöntemin ana yapılandırmaları 5 tohumla tekrarlanmıştır. Diğer görüşler, gerçekçilik testi, kök neden analizi, savunma deneyleri, bozuk bağlam deneyi ve bütçe taraması tek tohumla yapılmıştır.
\item \textbf{Savunma saldırganları.} Savunma saldırganları daha az veriyle eğitildiği için raporlanan sızıntı alt sınır niteliğindedir.
\item \textbf{Açıklanabilirlik analizi.} Grad-CAM haritaları $7 \times 7$ çözünürlükten büyütüldüğü için bölge payları yön göstericidir. Kesit yönü tahmini sagital kesitlerde güvenilir, aksiyel ve koronal kesitlerde kısmen doğrudur.
\item \textbf{Savunma ayrıntısı.} Sızıntı-güdümlü seçim $8 \times 8$ yama ızgarasıyla sınırlıdır; daha ince bir ızgara biraz daha verimli seçimler bulabilir.
\item \textbf{Maliyet ölçümü.} Maliyet tek bir bilgisayarda, $\Pi_{\mathrm{ROI}}$'deki lineer hesaplama yapısıyla ve önerilen yöntemin sığ evrişimli modeliyle ölçülmüştür; derin evrişimli ağlar için şifreli çıkarım maliyeti ölçülmemiştir.
\item \textbf{Önerilen yöntemin varsayımları.} FoveaHE, ROI'nin istemcide bilindiğini varsayar. Yeni bir güvenlik kanıtı sunmaz; sızıntı denetimi ampiriktir ve zaman yan kanalı tek bir bilgisayarda ölçülmüştür. Kullanılan şifreli modeller sığdır ve tek bir CKKS parametre kümesi denenmiştir.
\item \textbf{Rakip yöntem ve harici test.} Şifreli özet yaklaşımı doğrulukta FoveaHE'den üstündür. Harici genelleme testi tamamen bağımsız bir hastane verisi değildir ve harici küme iç kümeden daha kolaydır.
\item \textbf{Tehdit modeli.} Yalnızca yarı dürüst sunucu incelenmiştir.
\end{itemize}
